% from ega3-6-fr.tex
\oldpage[III]{144}
Soit $\rho:A\to A'$ un homomorphisme d'anneaux. Nous nous proposons de
définir un $A$-homomorphisme factoriel de degré 0 canoniquement associé à
$\rho$ :
\[
  \rho_{P_\bullet,Q_\bullet}:
  \mathbf{Tor}_\bullet^A(P_\bullet,Q_\bullet)
  \longrightarrow
  \mathbf{Tor}_\bullet^{A'}
  (P_\bullet\otimes_A A',Q_\bullet\otimes_A A').
  \tag{6.3.8.1}\label{III.6.3.8.1.fr}
\]

Pour cela, considérons une résolution de Cartan-Eilenberg projective
$L_{\bullet\bullet}$ de $P_\bullet$; considérons d'autre part une résolution
de Cartan-Eilenberg projective $L'_{\bullet\bullet}$ de
$P_\bullet\otimes_A A'$. Nous allons voir qu'on peut définir un
$A'$-homomorphisme de complexes
$L_{\bullet\bullet}\otimes_A A'\to L'_{\bullet\bullet}$, déterminé à
homotopie près. En effet, la construction de $L_{\bullet\bullet}$ est
entièrement déterminée lorsqu'on se donne (arbitrairement) pour chaque $i$,
une résolution projective $(X_{ij}^{\mathrm B})_{j\geqslant0}$ de
$\mathrm B_i(P_\bullet)$ et une résolution projective
$(X_{ij}^{\mathrm H})_{j\geqslant0}$ de $\mathrm H_i(P_\bullet)$, qui sont
respectivement égales à $\mathrm B_i^{\mathrm I}(L_{\bullet\bullet})$ et
$\mathrm H_i^{\mathrm I}(L_{\bullet\bullet})$; on en déduit successivement
\[
  \mathrm Z_i^{\mathrm I}(L_{\bullet\bullet})
  =\mathrm H_i^{\mathrm I}(L_{\bullet\bullet})
   \oplus\mathrm B_i^{\mathrm I}(L_{\bullet\bullet}),
  \qquad
  L_{i,\bullet}=\mathrm Z_i^{\mathrm I}(L_{\bullet\bullet})
   \oplus\mathrm B_{i-1}^{\mathrm I}(L_{\bullet\bullet}).
\]
Cela étant, $X_{i,\bullet}^{\mathrm B}\otimes_A A'$ n'est plus en général
une résolution de $\mathrm B_i(P_\bullet)\otimes_A A'$, mais est encore un
complexe formé de $A'$-modules projectifs, et il y a donc un
$A'$-homomorphisme
\[
  X_{i,\bullet}^{\mathrm B}\otimes_A A'
  \longrightarrow \mathrm B_i^{\mathrm I}(L'_{\bullet\bullet})
\]
compatible avec les augmentations, et déterminé à homotopie près
(M, V, 1.1). On a de même un $A'$-homomorphisme
\[
  X_{i,\bullet}^{\mathrm H}\otimes_A A'
  \longrightarrow \mathrm H_i^{\mathrm I}(L'_{\bullet\bullet})
\]
déterminé à homotopie près, d'où l'on déduit, par la construction rappelée
plus haut, un $A'$-homomorphisme
$L_{i,\bullet}\otimes_A A'\to L'_{i,\bullet}$ pour tout $i$; ces
homomorphismes (pour $i\in\mathbf Z$) sont compatibles avec les opérateurs de
dérivation $L_{i,\bullet}\to L_{i-1,\bullet}$ et les analogues pour
$L'_{\bullet\bullet}$, en vertu de la même construction, et ils constituent
donc le $A'$-homomorphisme
$L_{\bullet\bullet}\otimes_A A'\to L'_{\bullet\bullet}$ cherché.

Pour définir (6.3.8.1), il suffit alors de considérer de même une résolution
de Cartan-Eilenberg projective $M_{\bullet\bullet}$ (resp.
$M'_{\bullet\bullet}$) de $Q_\bullet$ (resp.
$Q_\bullet\otimes_A A'$), et un $A'$-homomorphisme
$M_{\bullet\bullet}\otimes_A A'\to M'_{\bullet\bullet}$. On déduit de ces
homomorphismes un $A'$-homomorphisme
\[
  (L_{\bullet\bullet}\otimes_A A')\otimes_{A'}
  (M_{\bullet\bullet}\otimes_A A')
  \longrightarrow
  L'_{\bullet\bullet}\otimes_{A'}M'_{\bullet\bullet},
\]
puis par composition un $A$-homomorphisme de bicomplexes
\[
  L_{\bullet\bullet}\otimes_A M_{\bullet\bullet}
  \longrightarrow
  L'_{\bullet\bullet}\otimes_{A'}M'_{\bullet\bullet},
\]
et en passant à l'homologie on obtient (6.3.8.1), qui est bien défini
puisqu'il provient d'un morphisme de complexes défini à homotopie près.

Si $\rho':A'\to A''$ est un second homomorphisme d'anneaux, et
$\rho'':A\to A''$ l'homomorphisme composé $\rho'\circ\rho$, il est clair que
\[
  \rho''_{P_\bullet,Q_\bullet}
  =\rho'_{P'_\bullet,Q'_\bullet}\circ\rho_{P_\bullet,Q_\bullet},
  \qquad
  P'_\bullet=P_\bullet\otimes_A A',
  \quad Q'_\bullet=Q_\bullet\otimes_A A'.
\]
Notons encore que le morphisme de bicomplexes
$L_{\bullet\bullet}\otimes_A M_{\bullet\bullet}
\to L'_{\bullet\bullet}\otimes_{A'}M'_{\bullet\bullet}$ considéré
ci-dessus définit des morphismes factoriels (en $P_\bullet$ et $Q_\bullet$)
de suites spectrales
\[
  'E_{pq}^r(P_\bullet,Q_\bullet)
  \longrightarrow
  'E_{pq}^r(P_\bullet\otimes_A A',Q_\bullet\otimes_A A')
\]
et
\[
  ''E_{pq}^r(P_\bullet,Q_\bullet)
  \longrightarrow
  ''E_{pq}^r(P_\bullet\otimes_A A',Q_\bullet\otimes_A A'),
\]
indépendants des résolutions de Cartan-Eilenberg considérées, et ayant aussi
la propriété de transitivité précédente.
\end{env}

\begin{proposition}[6.3.9]
Soit $\rho:A\to A'$ un homomorphisme d'anneaux tel que $A'$ soit un
$A$-module plat. On a alors des isomorphismes canoniques factoriels
\[
  \mathbf{Tor}_\bullet^{A'}
  (P_\bullet\otimes_A A',Q_\bullet\otimes_A A')
  \xrightarrow{\sim}
  \mathbf{Tor}_\bullet^A(P_\bullet,Q_\bullet)\otimes_A A'.
  \tag{6.3.9.1}\label{III.6.3.9.1.fr}
\]
\[
  \left\{
  \begin{aligned}
  'E(P_\bullet\otimes_A A',Q_\bullet\otimes_A A')
    &\xrightarrow{\sim}'E(P_\bullet,Q_\bullet)\otimes_A A',\\
  ''E(P_\bullet\otimes_A A',Q_\bullet\otimes_A A')
    &\xrightarrow{\sim}''E(P_\bullet,Q_\bullet)\otimes_A A'.
  \end{aligned}
  \right.
  \tag{6.3.9.2}\label{III.6.3.9.2.fr}
\]
\end{proposition}

\begin{proof}
